\paragraph{}Este capítulo apresenta a fundamentação teórica essencial para o entendimento das metodologias e soluções computacionais propostas neste trabalho. A \autoref{sec:def_problema} apresenta o conceito do Problema de Roteamento de Veículos, detalha os padrões de avaliação de instâncias da literatura e formaliza-o matematicamente. A \autoref{sec:selecao_meta} apresenta as meta-heurísticas adotadas (Recozimento Simulado, Busca Tabu e Algoritmo Genético) descrevendo sucintamente seus funcionamentos e premissas. Posteriormente, a \autoref{sec:sma} explora os princípios de um Sistema multiagente, abordando o \textit{framework} de modelagem distribuída escolhido e a estratégia de cooperação implementada, o \textit{Path Relinking}. A compreensão do Aprendizado por Reforço e a lógica da recompensa empregada na estratégia do \textit{Q-Learning}, que norteia a coordenação estratégica proposta, são objetos da \autoref{sec:rl}. Por fim, a \autoref{sec:avaliacao_estatistica} discute a importância de uma avaliação estatística rigorosa para a comparação de algoritmos meta-heurísticos, destacando os testes de hipóteses adequados para validar as conclusões obtidas a partir dos experimentos computacionais realizados.

\section{Definição do problema}
\label{sec:def_problema}

\paragraph{}Para compreender a avaliação das soluções quantitativas na experimentação futura, faz-se necessário contextualizar a natureza exata e as particularidades do problema central que pauta toda a sequência teórica e empírica deste estudo. Inicialmente, a \autoref{sec:vrp} define o arranjo lógico do conhecido e amplamente mapeado VRP, englobando as dificuldades computacionais e as raízes atadas à complexidade tempo de polinomial não-determinístico difícil (\textit{NP-hard}). Secundariamente, a \autoref{sec:benchmark} elucida a base paramétrica de referências extraída do Benchmark de Solomon de 1987, estabelecendo um padrão ouro viável de contraste na apuração da acurácia e otimização obtida pelas composições desenvolvidas em tempo de execução comparado.

\subsection{O Problema de Roteamento de Veículos}
\label{sec:vrp}

\paragraph{} O Problema de Roteamento de Veículos (VRP) é um desafio de otimização combinatória e programação linear inteira que busca determinar o conjunto ideal de rotas para uma frota de veículos entregarem mercadorias a um conjunto específico de clientes \cite{Dantzig59, Vidal20}. De forma genérica, o objetivo fundamental do VRP é minimizar o custo total da operação (geralmente expresso em distância total ou tempo) garantindo que:

\begin{itemize}
    \item Cada cliente deve ser visitado exatamente uma vez;
    \item Todos os veículos devem começar e terminar suas rotas em um depósito central;
    \item As restrições operacionais (como capacidade de carga ou janelas de tempo) devem ser estritamente respeitadas.
\end{itemize}

\paragraph{} O VRP é classificado na teoria da complexidade computacional como um problema \textit{NP-hard} (\textit{Non-deterministic Polynomial-time hard}) \cite{Garey79}. Do ponto de vista teórico, isso implica que qualquer problema pertencente à classe de complexidade \textit{NP} (\textit{Non-deterministic Polynomial}) pode ser reduzido ao VRP por meio de uma transformação em tempo polinomial. Na prática, isso significa que, a menos que $P = NP$, não se conhece nenhum algoritmo determinístico capaz de encontrar a solução ótima para o problema em tempo polinomial. É precisamente essa complexidade computacional inerente que fundamenta a aplicação das meta-heurísticas (SA, TS e GA) investigadas neste estudo. Em contraste com os métodos de resolução exatos, as abordagens meta-heurísticas são concebidas para encontrar soluções subótimas de alta qualidade em um tempo de processamento que viabilize a aplicação prática em cenários de envergadura real.

\subsection{Metodologia de avaliação e o Benchmark de Solomon}
\label{sec:benchmark}

\paragraph{}Para assegurar a validade científica deste estudo e permitir a comparabilidade dos resultados obtidos pelas meta-heurísticas propostas, a modelagem será conduzida estritamente sob os parâmetros do Benchmark de Solomon (1987). A utilização de instâncias consagradas na literatura é imperativa, pois fornecem um ``gabarito'' técnico que permite atestar a precisão e a eficiência de cada algoritmo desenvolvido.

\paragraph{}Dado que as instâncias de Solomon incorporam coordenadas geográficas, demandas de carga e janelas temporais, o problema aqui tratado é classificado como Problema de Roteamento de Veículos com Janelas de Tempo (VRPTW). Embora simplificada em relação às complexidades imprevisíveis do mundo real (como variações de tráfego ou clima), essa modelagem oferece um ambiente controlado e puramente matemático para a análise de desempenho algorítmico. Para que a comparação seja viável e metodologicamente aceitável, o instanciação seguirá as métricas rigorosas do \textit{benchmark}:

\begin{enumerate}
    \item \textbf{Função Objetivo Hierárquica:} A avaliação do desempenho seguirá uma prioridade de dois níveis:
    \begin{enumerate}
        \item \textbf{Minimização do número de veículos ($K$):} O objetivo primário é atender a todos os clientes com a menor frota possível.
        \item \textbf{Minimização da distância total percorrida ($D$):} Como objetivo secundário, busca-se o percurso mais curto para a frota definida.
    \end{enumerate}

    \item \textbf{Cálculo de Distância e Tempo:} A distância entre dois pontos $(i, j)$ é calculada pela Distância Euclidiana plana:
    
    $$d_{ij} = \sqrt{(x_i - x_j)^2 + (y_i - y_j)^2}$$

    \item Conforme as regras de Solomon, o tempo de viagem é numericamente equivalente à distância ($t_{ij} = d_{ij}$), e todos os cálculos devem utilizar precisão dupla (\textit{double precision}), com os resultados finais de distância arredondados para duas casas decimais.

    \item \textbf{Ponto de Origem e Retorno:} O cliente indexado como 0 define o depósito central, ponto obrigatório de partida e encerramento de todas as rotas.
    
    \item \textbf{Restrições de Viabilidade:} Uma solução para o problema é considerada válida apenas se garantir o atendimento da demanda de todos os clientes em sua totalidade. Além disso, cada rota individual deve respeitar rigorosamente as seguintes restrições operacionais: a soma das demandas dos clientes inseridos em uma mesma rota não pode exceder a Capacidade Máxima do Veículo ($Q$); e o serviço deve ser iniciado e finalizado dentro da Janela de Tempo estabelecida pelo cliente (\textit{Ready Time} e \textit{Due Date}), computando-se obrigatoriamente o Tempo de Serviço (\textit{Service Time}) em cada parada.
\end{enumerate}

\paragraph{}Esta abordagem padronizada transforma o problema em um estudo de caráter matemático e exploratório, estabelecendo uma base sólida que pode, futuramente, ser adaptada para modelos logísticos de maior complexidade e variáveis dinâmicas.

\subsection{Formulação Matemática}
\label{sec:modelo_milp}

\paragraph{} Para formalizar o VRPTW e estabelecer o modelo de referência deste trabalho, apresenta-se a seguir uma formulação de Programação Linear Inteira Mista (MILP -- \textit{Mixed-Integer Linear Programming}).

\paragraph{} Considere um grafo direcionado $G = (V, A)$, onde $V = \{0, 1, \dots, n, n+1\}$ é o conjunto de vértices e $A$ representa o conjunto de arcos válidos, livre de laços conexos ao próprio nó (\textit{self-loops}). Os vértices $0$ e $n+1$ delimitam geometricamente o depósito central em suas instâncias de partida e chegada, respectivamente, enquanto o subconjunto $C = \{1, \dots, n\}$ define os nós de demanda associados aos clientes. Para cada cliente $i \in C$, estipula-se uma demanda de carga $q_i$, um tempo de serviço fixo $s_i$ e uma janela temporal de atendimento restrita ao intervalo $[e_i, l_i]$. Cada arco $(i, j) \in A$ carrega um custo de percurso $c_{ij}$ e um tempo de deslocamento $t_{ij}$. 

\paragraph{} Para garantir que o tamanho da frota disponível $|K|$ atue estritamente como uma cota superior e não como uma imposição de despacho mandatória, o conjunto de arcos $A$ inclui explicitamente o arco direcionado virtual $(0, n+1)$, cujos parâmetros são zerados ($c_{0, n+1} = 0$ e $t_{0, n+1} = 0$). Desse modo, qualquer veículo homogêneo $k \in K$ que permaneça ocioso cumpre os requisitos de partida e chegada realizando a trajetória direta entre os depósitos virtuais sem penalizar a função objetivo. 

\paragraph{} O modelo adota duas classes de variáveis de decisão: a variável binária $x_{ijk}$, que assume valor 1 se o veículo $k \in K$ percorre diretamente o arco $(i, j) \in A$, e 0 caso contrário; e a variável contínua $w_i \ge 0$, que estipula o instante de início do serviço no vértice $i \in V$. A formulação matemática exata é estruturada conforme segue:

\begin{align}
\text{Minimizar} \quad & Z = \sum_{k \in K} \sum_{(i,j) \in A} c_{ij}\, x_{ijk}
  \label{eq:obj}\\[4pt]
\text{Sujeito a:} \quad
& \sum_{k \in K} \sum_{j:(i,j)\in A} x_{ijk} = 1,
  && \forall i \in C \label{eq:grau}\\[2pt]
& \sum_{j:(0,j)\in A} x_{0jk} = 1,
  && \forall k \in K \label{eq:saida}\\[2pt]
& \sum_{i:(i,h)\in A} x_{ihk} - \sum_{j:(h,j)\in A} x_{hjk} = 0,
  && \forall h \in C,\ \forall k \in K \label{eq:fluxo}\\[2pt]
& \sum_{i:(i,n+1)\in A} x_{i,n+1,k} = 1,
  && \forall k \in K \label{eq:chegada}\\[2pt]
& \sum_{i \in C} q_i \sum_{j:(i,j)\in A} x_{ijk} \le Q,
  && \forall k \in K \label{eq:cap}\\[2pt]
& w_i + s_i + t_{ij} - M_{ij}\bigl(1 - \textstyle\sum_{k\in K} x_{ijk}\bigr) \le w_j,
  && \forall (i,j) \in A \label{eq:tempo}\\[2pt]
& e_i \le w_i \le l_i,
  && \forall i \in V \label{eq:janela}\\[2pt]
& x_{ijk} \in \{0,1\},
  && \forall (i,j) \in A,\ \forall k \in K \label{eq:dominio}
\end{align}

\paragraph{} A função objetivo \eqref{eq:obj} visa minimizar o custo total de deslocamento da frota. A restrição de grau \eqref{eq:grau} impõe que cada nó de cliente seja visitado exatamente uma vez por um único veículo. As restrições \eqref{eq:saida}, \eqref{eq:fluxo} e \eqref{eq:chegada} gerenciam as leis de conservação de fluxo, obrigando que todo veículo ativado saia do depósito inicial $0$, preserve a continuidade física ao transitar por nós intermediários de clientes e finalize seu percurso no depósito de encerramento $n+1$. 

\paragraph{} Os limites físicos de carregamento de cada veículo são controlados pela inequação \eqref{eq:cap}. A coordenação cronológica e a eliminação de subrotas disjuntas (\textit{subtour elimination}) são coordenadas pelas expressões \eqref{eq:tempo} e \eqref{eq:janela}. A fim de prover um relaxamento linear computacionalmente mais restritivo e robusto se comparado ao uso de uma constante global arbitrária, a Equação \eqref{eq:tempo} adota um coeficiente penalizador Big-$M$ indexado localmente, parametrizado por $M_{ij} = \max\{0, l_i + s_i + t_{ij} - e_j\}$. Por fim, a restrição \eqref{eq:janela} delimita as fronteiras operacionais das janelas de tempo de cada vértice.

\paragraph{Relação com a Abordagem Heurística}

\paragraph{} Determinar a solução ótima do modelo MILP por meio de métodos exatos é uma tarefa intratável para problemas de grande porte devido à classificação do VRPTW como um problema \textit{NP-Hard}. Frente a essa limitação, este trabalho recorre a abordagens meta-heurísticas para guiar a busca por soluções de alta qualidade.

\paragraph{} Sob a perspectiva algorítmica, o espaço de busca e a validação das restrições matemáticas sofrem uma reconfiguração metodológica, dividindo-se em dois pilares:
\begin{enumerate}
    \item \textbf{Atendidas por construção via representação:} Ao contrário dos solucionadores exatos, que avaliam tensores de adjacência infactíveis durante o processo de ramificação, os algoritmos heurísticos propostos eliminam a necessidade de lidar com inviabilidades de rota. As restrições de conservação de fluxo \eqref{eq:saida}-\eqref{eq:chegada}, capacidade \eqref{eq:cap}, sincronismo temporal \eqref{eq:tempo} e janelas de atendimento \eqref{eq:janela} são asseguradas de forma intrínseca pela estrutura de codificação adotada, impossibilitando o software de gerar um plano de transporte operacionalmente inválido.
    \item \textbf{Aproximadas empiricamente:} A otimização estrita da função objetivo \eqref{eq:obj} deixa de ser avaliada sob uma perspectiva global fechada. O custo total e o dimensionamento da frota passam a ser minimizados de forma empírica por meio da dinâmica cooperativa e competitiva dos operadores estocásticos, os quais sacrificam a prova matemática de otimalidade em prol da viabilidade de tempo de execução computacional.
\end{enumerate}


\subsection{Representação de soluções: O paradigma de \textit{giant tour}}
\label{subsec:giant_tour}

\paragraph{} O fator determinante para o desempenho de meta-heurísticas aplicadas ao VRPTW reside na escolha da estrutura de dados empregada para codificar uma solução. A manipulação direta da álgebra linear exata apresentada na \autoref{sec:modelo_milp} por operadores estocásticos impõe barreiras geométricas intransponíveis à convergência dos algoritmos. Para contornar essa restrição estrutural, a literatura de Pesquisa Operacional frequentemente adota o paradigma de representação baseado em rotas gigantes, ou \textit{giant tours}.

\paragraph{} Proposto por \cite{prins2004simple}, o conceito de \textit{giant tour} consiste em transpor uma solução complexa multirrota do VRP para uma única sequência linear e contínua contendo a totalidade dos $n$ clientes mapeados, omitindo as inserções intermediárias do depósito central. Essa abstração permite que o problema de roteamento seja operado temporariamente sob as propriedades de uma permutação pura do Problema do Caixeiro Viajante. Consequentemente, as meta-heurísticas concentram seus mecanismos de busca exclusivamente na identificação da sequência de atendimento global ideal, explorando um espaço amostral significativamente mais denso e homogêneo.

\paragraph{} A restauração da viabilidade logística e a conversão do arranjo unidimensional em um plano de expedição multifrota ocorrem por meio de um decodificador de partição, denominado fatiamento ou \textit{split}. Neste projeto, codificou-se uma variante estabelecida sob o princípio da divisão gulosa (\textit{greedy split}), cuja execução ocorre de forma sequencial adjacente ao \textit{giant tour}. O algoritmo realiza a varredura progressiva da lista, acumulando as demandas de carga e os tempos de chegada nos nós. No instante exato em que a transição para o cliente subsequente acusa a iminência de violação dos limites de capacidade do veículo \eqref{eq:cap} ou estouramento do limite superior da janela temporal \eqref{eq:janela}, o decodificador força dinamicamente o retorno do veículo atual ao depósito central e despacha uma nova unidade para dar continuidade ao atendimento.

\paragraph{} Posto que o retorno ao depósito é desencadeado de forma reativa e automática diante de uma saturação operacional, evidencia-se que a totalidade das rotas fatiadas é viável por construção, eliminando por completo a ocorrência de soluções logísticas inválidas ao longo dos passos do algoritmo. No entanto, essa simplificação estrutural introduz um compromisso metodológico (\textit{trade-off}) que deve ser formalizado como uma limitação assumida do modelo. Devido ao caráter míope inerente à partição gulosa, que efetua os cortes de rota de maneira imediata sem mensurar o impacto combinatório dessa quebra nas demandas remanescentes, o procedimento tende a superestimar o número total de veículos ($K$) necessários para consolidar a frota final. Trata-se de uma fragilidade conhecida, tolerada institucionalmente em favor da robustez estrutural e da velocidade de convergência computacional das meta-heurísticas em problemas de larga escala.

\section{Meta-heurísticas selecionadas}
\label{sec:selecao_meta}

\paragraph{}As meta-heurísticas Recozimento Simulado, Busca Tabu e Algoritmo Genético são clássicas no contexto do VRP, portanto adequadas para o presente estudo. A diversidade técnica entre esses métodos permite utilizar diferentes mecanismos de escape de ótimos locais e estratégias de exploração do espaço de busca diante das restrições de roteamento. A seguir, serão apresentados os fundamentos do funcionamento de cada um deles.

\subsection{Recozimento Simulado}
\label{sec:sa}

\paragraph{} Explora possíveis soluções por meio de um mecanismo probabilístico inspirado no processo de resfriamento de metais \cite{Kirkpatrick83}. Se $\Delta f$ é a diferença de custo entre a nova solução ($s'$) e a atual ($s$):

$$\Delta f = f(s') - f(s)$$

\vspace{0.5cm}
\noindent a probabilidade $P$ de aceitação é definida por:

$$P(\text{aceitar } s') = 
\begin{cases} 
1, & \text{se } \Delta f < 0 \\
e^{-\frac{\Delta f}{T}}, & \text{se } \Delta f \geq 0 
\end{cases}$$

\vspace{0.5cm}
\noindent onde $T$ é a temperatura, que decai ao longo do tempo conforme uma taxa de resfriamento $\alpha$:

$$T_{k+1} = \alpha \cdot T_k$$

\subsection{Busca Tabu}
\label{sec:ts}

\paragraph{} O método utiliza ``listas tabu'', que impedem o algoritmo de retornar a soluções visitadas recentemente, forçando uma exploração sistemática e agressiva de novas vizinhanças \cite{Glover89}. Mais objetivamente, o algoritmo opera sobre a vizinhança $N(s)$ da solução atual, mas subtrai os movimentos que estão na Lista Tabu ($L$):

$$s_{next} = \text{argmin} \{f(s') \mid s' \in N(s) \setminus L\}$$

\paragraph{}Entretanto, existe o Critério de Aspiração. Se uma solução $s'$ está na lista tabu, mas seu custo é melhor que a melhor solução encontrada até agora ($s_{best}$), o estado tabu é ignorado:

$$\text{Se } f(s') < f(s_{best}), \text{ então } s' \in \textit{Candidatos} \text{ mesmo se } s' \in L$$

\subsection{Algoritmo Genético}
\label{sec:ga}

\paragraph{} Por meio dos operadores de cruzamento, mutação e seleção, o algoritmo é capaz de recombinar sub-rotas eficientes de diferentes soluções para gerar descendentes potencialmente superiores \cite{Holland75, Goldberg89}. Isso se traduz em converter o problema de minimização de custos em um valor de sobrevivência, representado pela função de aptidão (\textit{fitness}) $F(s)$:

$$F(s) = \frac{1}{C(s) + \epsilon}$$

\vspace{0.5cm}
\noindent onde $\epsilon$ é um valor infinitesimal para evitar divisão por zero. Presente, neste caso, apenas por rigor teórico, pois a função de custo que será utilizada é sempre maior que 0. A probabilidade de um indivíduo ser selecionado para o cruzamento é dada por:

$$P_i = \frac{F(s_i)}{\sum_{j=1}^{N} F(s_j)}$$

\vspace{0.5cm}
\noindent onde $N$ é o tamanho da população. Isso garante que rotas mais curtas tenham maior probabilidade biológica de transmitir seus genes (sequências de clientes).


\section{Sistema multiagente}
\label{sec:sma}

\paragraph{}Os \textit{Multi-Agent Systems} (MAS) consistem em um conjunto de elementos computacionais, denominados agentes, que interagem em um ambiente compartilhado para resolver problemas que ultrapassam as capacidades individuais de um único processador ou algoritmo \cite{Wooldridge95}. Tais sistemas têm sido aplicados com sucesso em diversas áreas, incluindo otimização de colônias de formigas, onde a emergência de comportamentos complexos globais decorre das interações locais entre agentes relativamente simples.

\paragraph{}No contexto do VRPTW, a abordagem multiagente é empregada para decompor a complexidade da busca em diferentes especialidades: enquanto um agente foca na exploração global do espaço de busca (\textit{exploration}), outro dedica-se ao refinamento local intensivo (\textit{exploitation}).

\paragraph{}Essa arquitetura permite a transição de algoritmos competitivos para um modelo cooperativo, onde a inteligência coletiva emerge da troca de informações entre os agentes. A principal vantagem desta abordagem para o VRPTW é a capacidade de evitar que o sistema fique preso em ótimos locais, uma vez que diferentes comportamentos heurísticos operam simultaneamente e compartilham suas melhores descobertas por meio de uma memória comum.

\subsection{Framework Mesa}

\paragraph{}O Mesa é um framework de código aberto desenvolvido em \textit{Python} para a modelagem baseada em agentes, permitindo que pesquisadores criem, manipulem e analisem sistemas complexos compostos por múltiplos elementos autônomos \cite{TerHoeven25}. Diferente de modelos matemáticos puramente estáticos, o Mesa fornece uma estrutura computacional que facilita a observação de comportamentos emergentes a partir das interações locais entre agentes em um ambiente dinâmico.

\paragraph{}A arquitetura do Mesa é fundamentada em uma separação clara de responsabilidades, baseada em três componentes principais que seguem padrões de projeto modernos:

\begin{itemize}
    \item \textbf{Agente (\textit{Agent}):} Representa a unidade fundamental de inteligência e ação. Cada agente possui seu próprio estado e lógica interna de decisão, sendo capaz de interagir com o ambiente ou com outros agentes.
    
    \item \textbf{Modelo (\textit{Model}):} Atua como a entidade de alto nível que orquestra o sistema. O modelo define as regras do ambiente, mantém os dados globais e gerencia a criação e destruição dos agentes.
    
    \item \textbf{Agendador (\textit{Scheduler}):} É o componente responsável pela coordenação temporal das ações. Ele determina a ordem em que os agentes executam seus passos (\textit{steps}), podendo operar de forma aleatória, sequencial ou simultânea, o que impacta diretamente a dinâmica de convergência do sistema.
\end{itemize}

No âmbito da otimização combinatória, o Mesa destaca-se por permitir que meta-heurísticas distintas sejam encapsuladas em agentes independentes, transformando um problema de busca individual em um processo de cooperação distribuída.

\subsection{Estratégia de cooperação: \textit{Path Relinking}}

\paragraph{} O \textit{Path Relinking} (PR) é uma estratégia de intensificação de busca originalmente proposta como um componente de memória para a Busca Tabu \cite{Glover00}, mas que se consolidou como uma técnica poderosa para a integração de soluções em diversos contextos de otimização combinatória.

\paragraph{} O conceito fundamental do PR reside na exploração sistemática de trajetórias que conectam duas soluções de alta qualidade, denominadas solução inicial ($S_i$) e solução guia ($S_g$). A premissa teórica do método é que, ao navegar pelo caminho que transforma a solução inicial na solução guia, o algoritmo pode descobrir soluções intermediárias que compartilham características de ambas e que, em muitos casos, superam as referências originais em termos de qualidade.

\paragraph{} Esse processo é realizado por meio de uma sequência de movimentos de vizinhança que reduzem progressivamente a distância entre as duas soluções. No escopo deste trabalho, a distância de Hamming é calculada diretamente sobre a representação do vetor de permutação linear dos clientes (a estrutura do \textit{giant tour}). Registra-se, contudo, como limitação metodológica assumida, que essa métrica afere estritamente a divergência posicional dos elementos nos índices do arranjo numérico; logo, ela não é capaz de capturar a estrutura de arcos logísticos conetivos ou o compartilhamento de arestas reais entre as rotas físicas do grafo. Embora o conceito fundamental descreva uma trajetória estrutural entre dois pontos, a literatura clássica \cite{Glover00} categoriza o \textit{Path Relinking} em diferentes tipos, dependendo de como e em qual direção essa trajetória é explorada:

\begin{itemize}
    \item \textbf{\textit{Forward Path Relinking} (Direto):} A busca parte da pior solução (Solução Inicial - $S_i$) em direção à melhor solução (Solução Guia - $S_g$). A premissa é incorporar progressivamente as boas características da solução guia na solução inicial.
    
    \item \textbf{\textit{Backward Path Relinking} (Reverso):} A trajetória é percorrida no sentido inverso, partindo da melhor solução ($S_g$) em direção à pior ($S_i$). Esta variante costuma apresentar resultados superiores na prática, pois o algoritmo explora a vizinhança imediata da solução de mais alta qualidade logo nos primeiros passos da trajetória.
    
    \item \textbf{\textit{Mixed Path Relinking} (Misto):} A exploração ocorre de forma alternada, avançando simultaneamente a partir de ambas as extremidades ($S_i$ e $S_g$) até que os caminhos se encontrem no centro da trajetória, garantindo uma amostragem equilibrada de ambas as vizinhanças.
\end{itemize}

\paragraph{} Para o escopo do modelo de cooperação proposto neste trabalho, a implementação 
foca especificamente nas estratégias \textit{Forward} e \textit{Backward}, 
permitindo que os agentes explorem essas trajetórias bidirecionalmente para 
maximizar a extração de características promissoras das rotas.

\begin{figure}[H]
    \centering
    \parbox[htb]{13.0cm}
    {
        \centering
        \includegraphics[scale=0.5]{path_relinking.png}
        \caption[\small{Diagrama de Path Relinking}]{\label{FigPR} \small{Representação do caminho entre $S_i$ e $S_g$.}}
    }
\end{figure}
\vspace{-0.5cm}

\paragraph{} No contexto do VRPTW, o \textit{Path Relinking} atua como um mecanismo de cooperação de alto nível ao permitir a sinergia estrutural, combinando a eficiência de sequenciamento de uma solução com a estrutura de particionamento de veículos de outra para criar rotas que nenhum agente isolado seria capaz de conceber. Além disso, a técnica possibilita a exploração de vizinhanças distantes, uma vez que, ao contrário de buscas locais tradicionais que exploram apenas o entorno imediato, o PR investiga uma região específica do espaço de busca delimitada por duas soluções de elite, focando o esforço computacional em áreas comprovadamente promissoras.

\paragraph{} Outro fator determinante é a sua capacidade de facilitar o escape de ótimos locais, pois a trajetória entre as soluções pode atravessar vales de soluções inviáveis ou de baixa qualidade para alcançar picos de aptidão mais elevados, agindo como um catalisador para a convergência global do sistema multiagente. Dessa forma, o PR não é apenas um método de busca, mas uma filosofia de cooperação que utiliza o conhecimento acumulado por diferentes agentes especialistas para sintetizar soluções superiores, mitigando as fraquezas individuais observadas em execuções isoladas. 

\paragraph{} Como seus operadores e manipuladores de soluções são relativamente simples, o PR é menos custoso computacionalmente que uma meta-heurística, seu esforço reside majoritariamente na verificação das funções de custo. Este é seu ponto forte, mas também ponto fraco, uma vez que depende de boas soluções como entrada para melhorá-las ainda mais. É evidente, portanto, que deve ser usada em combinação com outros algoritmos, idealmente ao final das etapas.

\section{Aprendizado por Reforço}
\label{sec:rl}

\paragraph{} O Aprendizado por Reforço (\textit{Reinforcement Learning} - RL) é uma área do aprendizado de máquina que se preocupa com a forma como agentes inteligentes devem tomar ações em um ambiente a fim de maximizar uma recompensa acumulada. Diferente do aprendizado supervisionado, que se baseia em um conjunto de dados rotulados, fundamenta-se na experiência direta do agente, que aprende por meio de um processo de tentativa e erro ao interagir com o ambiente \cite{Sutton18}.

\paragraph{} A estrutura básica de um problema de RL é frequentemente modelada como um Processo de Decisão de Markov, que envolve os seguintes componentes fundamentais:

\begin{itemize}
    \item \textbf{Agente (\textit{Agent}):} O tomador de decisão que observa o ambiente e executa ações.
    \item \textbf{Ambiente (\textit{Environment}):} O contexto externo com o qual o agente interage.
    \item \textbf{Estado (\textit{State}):} Uma representação da situação atual do ambiente em um determinado instante. No contexto de otimização, o estado pode refletir a qualidade da solução atual, como o número de veículos e a distância total.
    \item \textbf{Ação (\textit{Action}):} O conjunto de decisões possíveis que o agente pode tomar.
    \item \textbf{Recompensa (\textit{Reward}):} Um sinal numérico de \textit{feedback} enviado pelo ambiente após uma ação, que indica o sucesso ou fracasso daquela decisão em relação ao objetivo final.
\end{itemize}

\paragraph{} O objetivo do RL é encontrar uma política ($\pi$), que é um mapeamento de estados para ações que maximizem o retorno esperado a longo prazo. Um conceito crítico neste processo é o equilíbrio entre exploração (\textit{exploration}), onde o agente testa novas ações para descobrir sua utilidade, e explotação (\textit{exploitation}), onde o agente utiliza o conhecimento já adquirido para escolher as ações que renderam as maiores recompensas anteriormente.

\paragraph{} Um dos mais clássicos algoritmos de RL, O Q-Learning \cite{Watkins92} é do tipo \textit{off-policy} e \textit{model-free}, o que significa que ele pode aprender a política ótima de tomada de decisão sem a necessidade de um modelo prévio das transições do ambiente ou das probabilidades de recompensa. O ``Q'' na nomenclatura refere-se à Qualidade (ou valor de utilidade), representando o quão proveitosa é uma determinada ação em um estado específico.

\begin{figure}[H]
    \centering
    \begin{tikzpicture}[
        x=1cm,
        y=1cm,
        >=Stealth,
        frame/.style={dash pattern=on 2pt off 2pt, line width=0.4pt},
        box/.style={
            draw,
            rounded corners=2pt,
            line width=0.5pt,
            fill=gray!12,
            align=center,
            inner xsep=1.5pt,
            inner ysep=0.6pt,
            minimum height=0.9cm
        },
        flow/.style={-Stealth, line width=0.5pt}
    ]
        % Canvas and split between agent side (left) and environment side (right).
        \draw[frame] (-0.45,0) rectangle (5,7);

        \node[box, minimum width=3cm] (agent) at (2.5,5.5) {Agente de Atualização\\[-8pt]da Função-Q};
        \node[box, minimum width=3cm] (reward) at (2.5,3.5) {Avaliação de\\[-8pt]Recompensa/Penalidade};
        \node[box, minimum width=2cm] (actions) at (2.5,1.5) {Conjunto de\\[-8pt]Ações Possíveis};
        \node[box, minimum width=2cm] (obs) at (7.0,5.5) {Observação};
        \node[box, minimum width=2cm] (fav) at (7.0,1.5) {Ação Mais\\[-8pt]Favorável};
        \node[box, minimum width=2cm] (env) at (9.0,3.5) {Ambiente};

        \draw[flow] (obs.west) -- (agent.east);
        \draw[flow] (env.north) |- (obs.east);
        \draw[flow] (agent.south) -- (reward.north);
        \draw[flow] (reward.south) -- (actions.north);
        \draw[flow] ([xshift=0.45cm]actions.north) -- ([xshift=0.45cm]reward.south);
        \draw[flow] (actions.east) -- (fav.west);
        \draw[flow] (fav.east) -| (env.south);
        \coordinate (leftloop) at ([xshift=-0.2cm]reward.west);
        \draw[flow] (agent.west) -- (leftloop |- agent.west) -- (leftloop |- actions.west) -- (actions.west);

    \end{tikzpicture}
    \caption[\small{Visão geral do Q-Learning}]{\label{FigQL1} \small{Esquema de interação entre agente e ambiente no processo de aprendizagem.}}
\end{figure}
\vspace{-0.5cm}

\paragraph{} O núcleo do algoritmo consiste em estimar os valores da função $Q(s, a)$, que mensura a recompensa total esperada a longo prazo que um agente pode obter ao realizar a ação $a$ partindo do estado $s$. Em problemas com espaços de estados discretos, esses valores são armazenados em uma estrutura de dados conhecida como Tabela-Q (\textit{Q-Table}), onde as linhas representam os estados do sistema e as colunas as ações possíveis.

\begin{figure}[H]
    \centering
    \begin{tikzpicture}[
        >=Stealth,
        box/.style={
            draw,
            rounded corners=2pt,
            line width=0.5pt,
            fill=gray!12,
            align=center,
            inner xsep=2pt,
            inner ysep=1pt,
            minimum width=2.8cm,
            minimum height=1.0cm
        },
        flow/.style={-Stealth, line width=0.5pt}
    ]
        \node[box] (s1) at (0,0) {Inicializar\\[-8pt]Tabela-Q};
        \node[box] (s2) at (0,-1.8) {Escolher uma\\[-8pt]ação};
        \node[box] (s3) at (0,-3.6) {Executar\\[-8pt]ação};
        \node[box] (s4) at (0,-5.4) {Medir\\[-8pt]recompensa};
        \node[box] (s5) at (0,-7.2) {Atualizar\\[-8pt]Tabela-Q};

        \draw[flow] (s1.south) -- (s2.north);
        \draw[flow] (s2.south) -- (s3.north);
        \draw[flow] (s3.south) -- (s4.north);
        \draw[flow] (s4.south) -- (s5.north);
        \draw[flow] (s5.east) -- ++(1.9,0) -- ++(0,5.4) -- (s2.east);
        \node[font=\scriptsize, align=center, rotate=90] at (3.95,-4.5)
            {Após múltiplas iterações,\\[-1pt]uma boa Tabela-Q está pronta};
    \end{tikzpicture}
    \caption[\small{Algoritmo básico de Q-Learning}]{\label{FigQL2} \small{Diagrama do algoritmo de atualização do valor-$Q$.}}
\end{figure}
\vspace{-0.5cm} % Ajuste este valor conforme necessário para aproximar o texto

\paragraph{} O aprendizado ocorre através da atualização iterativa dos valores-Q (\textit{Q-values}) conforme o agente explora o ambiente. Essa atualização fundamenta-se na Equação de Bellman, que ajusta o valor atual com base na recompensa imediata recebida e na melhor estimativa de recompensas futuras:

\begin{equation}
    \label{EqQLearning}
    Q(s, a) \leftarrow Q(s, a) + \alpha [r + \gamma \max_{a'} Q(s', a') - Q(s, a)]
\end{equation}

\paragraph{} Nesta formulação, os parâmetros de controle desempenham papéis cruciais:

\begin{itemize}
    \item \textbf{Taxa de Aprendizado ($\alpha$):} Determina o quanto a nova informação substitui o conhecimento antigo.
    \item \textbf{Recompensa Imediata ($r$):} É o \textit{feedback} numérico direto recebido do ambiente após a execução de uma ação. Representa o ganho instantâneo e serve como a base factual para atualizar a estimativa de qualidade daquela ação.
    \item \textbf{Fator de Desconto ($\gamma$):} Define a importância das recompensas futuras em relação às imediatas.
    \item \textbf{Diferença Temporal:} O termo entre colchetes representa o erro de predição, ou seja, a diferença entre o valor que o agente esperava e o que ele efetivamente percebeu após a transição para o novo estado $s'$.
\end{itemize}

\paragraph{} Para que o algoritmo seja eficaz, o agente deve balancear a busca por novos conhecimentos com o uso do conhecimento já adquirido. A estratégia mais comum é a $\epsilon-greedy$, onde com uma probabilidade $\epsilon$ o agente escolhe uma ação aleatória (exploração) e com probabilidade $1 - \epsilon$ ele seleciona a ação com o maior valor-Q atual para aquele estado (explotação). Ao final do processo de treinamento, a Tabela-Q converge para uma representação da estratégia ótima, permitindo que o sistema, ao se deparar com uma determinada configuração do problema, saiba as melhores ações a serem tomadas.

\section{Avaliação estatística de desempenho algorítmico}
\label{sec:avaliacao_estatistica}

\paragraph{} A avaliação de desempenho de meta-heurísticas aplicadas a problemas de otimização combinatória, como o VRPTW, impõe desafios metodológicos devido à natureza estocástica inerente a esses algoritmos. A comparação direta de métricas de tendência central, como a média aritmética, é insuficiente para atestar a superioridade de um método sobre outro, uma vez que flutuações amostrais decorrentes de inicializações aleatórias (\textit{seeds}) podem distorcer os resultados. Torna-se imperativo, portanto, o emprego de testes de hipóteses estatísticas.

\paragraph{} Visto que as distribuições de custos de rotas e tempos de processamento computacional raramente atendem aos pressupostos clássicos de normalidade (distribuição Gaussiana) e homogeneidade de variância, a aplicação de testes paramétricos tradicionais, como o teste $t$ de Student ou a clássica ANOVA, torna-se inadequada. Em cenários de \textit{benchmarks} múltiplos avaliados sobre um mesmo conjunto de instâncias correlacionadas (amostras pareadas), a literatura prescreve a utilização de abordagens não-paramétricas baseadas na ordenação de resultados por postos (\textit{ranks}).

\subsection{Teste de Friedman}

\paragraph{} O Teste de Friedman configura-se como a alternativa não-paramétrica à ANOVA para medidas repetidas. Ele é utilizado para detectar diferenças globais significativas entre múltiplos tratamentos aplicados a diversos blocos de teste. No contexto computacional, os tratamentos correspondem aos diferentes algoritmos ou arquiteturas implementadas (tais como SA, TS, GA, MAS e MAS-RL), enquanto os blocos representam as instâncias de teste do \textit{benchmark} (como as instâncias de Solomon).

\paragraph{} O procedimento do teste consiste em ordenar os algoritmos individualmente dentro de cada instância, atribuindo o posto $1$ ao melhor desempenho (ex: menor número de veículos) e o posto $k$ ao pior (sendo $k$ o número total de algoritmos analisados). Em caso de empates numéricos, atribui-se a média dos postos em disputa a cada um dos algoritmos empatados. A hipótese nula ($H_0$) do teste estabelece que as amostras provêm da mesma população contínua, ou seja, assume que as distribuições dos postos de todos os métodos avaliados são estatisticamente equivalentes. Por sua vez, a hipótese alternativa ($H_1$) dita que existe uma diferença estatisticamente significativa na distribuição dos postos de, pelo menos, um dos algoritmos no grupo avaliado em relação aos demais.

\paragraph{} A estatística do teste ($\chi^2_F$) é computada com base na soma dos postos de cada algoritmo, avaliando a disparidade desses valores em relação ao cenário de equivalência estrita. Se o \textit{p-valor} associado gerado pela distribuição do teste for inferior ao nível de significância preestabelecido $\theta$ (tipicamente adotado como $0.05$ ou $0.10$), rejeita-se a hipótese nula.

\subsection{Teste Post-Hoc de Nemenyi}

\paragraph{} A rejeição da hipótese nula no Teste de Friedman estabelece apenas a existência de uma disparidade de desempenho no grupo geral de algoritmos (\textit{omnibus test}), falhando em especificar quais pares exatos de métodos diferem significativamente entre si. Para realizar tais comparações múltiplas par a par, mantendo o controle rigoroso sobre a inflação do Erro do Tipo I (falsos positivos), exige-se a aplicação de um teste \textit{post-hoc}. Neste escopo, o Teste de Nemenyi é a metodologia padrão recomendada para a área de aprendizado de máquina e otimização.

\paragraph{} O Teste de Nemenyi determina que o desempenho de dois algoritmos é considerado significativamente distinto se a diferença absoluta entre os seus \textit{ranks} médios superar um limiar matemático rigoroso denominado Diferença Crítica (CD, do inglês \textit{Critical Difference}). O valor numérico da CD é calculado de acordo com a equação:

$$\text{CD} = q_\alpha \sqrt{\frac{k(k + 1)}{6N}}$$

\paragraph{} Onde $q_\alpha$ representa o valor crítico tabelado baseado na distribuição de amplitude \textit{Studentized Range} para o nível de significância $\alpha$ (ou $\theta$) adotado; $k$ representa o número de algoritmos comparados no experimento; e $N$ consiste no número de instâncias ou blocos independentes avaliados.

\paragraph{} Graficamente, as conclusões deste teste \textit{post-hoc} são comumente sumarizadas e representadas pelo Diagrama de Diferença Crítica (\textit{Critical Difference Diagram}). Nele, os métodos computacionais são dispostos em um eixo horizontal de acordo com o seu respectivo \textit{rank} médio de desempenho. Algoritmos que não apresentam diferenças estatísticas significativas são agrupados por uma linha espessa de conexão, oferecendo uma percepção visual inequívoca acerca do ranqueamento global, das distâncias de performance e das reais equivalências operacionais entre os sistemas propostos.